% part: intuitionistic-logic
% chapter: propositions-as-types

\documentclass[../../../include/open-logic-chapter]{subfiles}

\begin{document}

\olchapter{int}{pty}{विधाने म्हणजे प्रकार}

\begin{editorial}
  करी--हॉवर्ड अनुरूपतेवरील या प्रकरणाचा हा \emph{अतिशय प्रायोगिक}
  मसुदा आहे. याला अधिक स्पष्टीकरण आणि प्रेरणा हवी; यात चुका
  व गाळलेले भागही असण्याची शक्यता आहे. सामान्यीकरणाची सिद्धता
  तपासून विस्तारित केली पाहिजे. गुणाकार-प्रकाराची उदाहरणे नाहीत.
  क्रमपरिवर्तन आणि सरलीकरण रूपांतरणांचा समावेश केलेला नाही.
  येथे मूलतः गृहीत धरलेल्या (प्रकारयुक्त) लॅम्डा कलनाविषयीही
  साहित्य आल्यावर हे प्रकरण अधिक समजेल. अत्यंत सावधगिरीने वापरा.
\end{editorial}

\olimport{introduction}
%\olimport{sequent-natural-deduction}
\olimport{proof-terms}
\olimport{proofs-to-terms}
\olimport{terms-to-proofs}
\olimport{types}
\olimport{reduction}
\olimport{type-preservation}
\olimport{normalization}

\OLEndChapterHook

\end{document}
